% !PRINTED 270
\weberpage{270}
We set
\weq{3}{\begin{aligned}
2A &= uv + (-1)^{\frac{n+1}{8}}(u_1v_1+u_2v_2),\\
B &= uvu_1v_1 + uvu_2v_2 + (-1)^{\frac{n+1}{8}}u_1v_1u_2v_2\\
&= \frac{2}{u_1v_1} + \frac{2}{u_2v_2} + (-1)^{\frac{n+1}{8}}\frac{2}{uv},
\end{aligned}}
so that the interchanges that belong together come out as
\weq{4}{\begin{array}{lll}
\omega, & A, & B,\\
-\dfrac{1}{\omega}, & \varrho\,A, & \varrho^2 B,\\
\omega+1, & e^{\frac{\pi i(n+1)}{12}}\sigma\,A, & e^{-\frac{\pi i(n+1)}{12}}\sigma^2 B.
\end{array}}

A product of the form $A^h B^k$ therefore takes on, through the two interchanges
\[
\Bigl(\omega, -\frac{1}{\omega}\Bigr), \quad (\omega, \omega+1)
\]
the factors
\[
\varrho^{h-k}, \quad e^{(h-k)\frac{(n+1)\pi i}{12}}\sigma^{h-k},
\]
which are $=1$ when $n+1$ is divisible by 3. We therefore now form, with numerical coefficients $M_{h,k}$, functions of the form
\weq{5}{\Phi_{a,c,\partial} = \Sigma\,M_{h,k}A^h B^k,}
in which, if $n \equiv 0$ or $\equiv 1$ (mod 3), $h$ and $k$ may take only such (integral) values whose differences $h-k$ leave the same residue on division by 3, while if $n \equiv -1$ (mod 3) they are not subject to this restriction. Such a function itself, or at least its third power, therefore satisfies a transformation equation of the first level. It remains finite, moreover, for all finite values of $j(\omega)$, and is consequently an \gesperrt{integral} algebraic function of $j(\omega)$. [The transformation of the second order $\begin{pmatrix} 1, & 1\\ -1, & 1 \end{pmatrix}$ need not be brought in here, since the totality of the $\Phi_{a,c,\partial}$ remains unchanged already under the substitution $(\omega, \omega+1)$, and not only under $(\omega, \omega+2)$. Cf. the remark at the beginning of §\,73.]

If, then, we determine the constants in $\Phi_{a,c,\partial}$ so that no negative powers of $q$ occur in the expansions of these functions, then all the $\Phi_{a,c,\partial}$ must be equal to one and the same constant.
